\scanpage{41}
\margin{\rule{0.35pt}{2.2\baselineskip}}d) Les formes quadratiques $q$ sur $B$ satisfaisant
\deleted{\unsure{les}}la condition
\[
 q(e_B)=1 \tag*{(63)}
\]

Les équivalences entre b), c), d) viennent
d'être explicitées, celle entre a) et b) a
été vu dans 5). [Il reste encore à expliciter
les relations directes entre a) et d), et rappeler
(à propos de l'équiv. a) et b)) la relation
\[
\Omega=\mathrm{Spec}(B) \tag*{(64)}
\]
lorsque $\Omega$ \inserted{$\subset F$} et une structure d'Algèbre sur $B$
se correspondent. Enfin, on aura
\[
q=\unsure{\mathcal N|\mathscr N|N}_{B/\mathcal O_S} \tag*{(65)}
\]

[[Pour vérifier ce dernier point, nous allons le
reformuler pour $n$ quelconque. Notons que
la donnée d'une section de $\Sigma^n(F/S)$, ou
encore d'une $n$-forme $f_n$ sur $B$ telle que
$f_n(e_B)=1$, ou encore d'un $n$-diviseur relatif
effectif $\Omega_n$ sur $S$, équivaut à la donnée d'une
$\mathcal O_S$-Algèbre \inserted{\unsure{commu.|comm.}} $C_n$ \inserted{loc. libre} de rang $n$, et d'un hom \inserted{surj}
\[
\underset{\substack{\wr\downarrow\\\underline{\mathrm{Sym}}^*(B)/(e_B-1)}}{\underline{\mathrm{Alg\,Aff}}(F)}\longrightarrow C_n
\]
ou ce qui revient au $=$, d'un hom de \inserted{$\mathcal O_S$-Mod.}
\[
 B\xrightarrow{g} C_n \tag*{(66)}
\]
% archive sheet number: 40
